\section{Functions with Environments: Closures}
\label{sec:closures}

\code{lecture5d.rkt} and \code{lecture5e.rkt} add \code{lambda} to the environment interpreter, back with an eager \code{let}. As the file puts it, ``once again we'll first do something that behaves `incorrectly' (or at least counter-intuitively) and then we'll switch to a fixed version'': the first attempt has dynamic scope, and the fix is the one of \Cref{sec:dynscope}.

\subsection{Attempt 1: A Function Is Its Parameter and Body}

\tyname{Expr} gets \code{lambdaC} and \code{appC} back, \code{parse} is the week~4 one (a list whose head is not a keyword is an application), and the result type is the \tyname{Value} of week~4, with function values written \code{lambdaV}. A binding holds a \tyname{Value}:
\begin{lstlisting}
(define-type Value
  [numV (n : number)]
  [lambdaV (x : symbol) (e : Expr)])

(define-type Binding
  [bind (name : symbol) (val : Value)])
;; Env, empty-env, extend-env, lookup: as in lecture5a, returning a Value

(define (eval-env (env : Env) (e : Expr)) : Value
  (type-case Expr e
    [numC (n) (numV n)]
    [lambdaC (x e) (lambdaV x e)]
    [plusC (e1 e2) (numV (+ (numV-n (eval-env env e1)) (numV-n (eval-env env e2))))]
    [timesC (e1 e2) (numV (* (numV-n (eval-env env e1)) (numV-n (eval-env env e2))))]
    [letC (x e1 e2)
          (let ([v1 (eval-env env e1)])
            (eval-env (extend-env (bind x v1) env) e2))]
    [appC (e1 e2)
          (let [(v1 (eval-env env e1))
                (v2 (eval-env env e2))]
            (eval-env (extend-env (bind (lambdaV-x v1) v2) env) (lambdaV-e v1)))]
    [idC (x) (lookup x env)]))

(define (eval (e : Expr))
  (eval-env empty-env e))
\end{lstlisting}
\code{+} and \code{*} unwrap their operands with \code{numV-n} directly, and \code{appC} takes the function apart with \code{lambdaV-x} and \code{lambdaV-e}, so adding a function or applying a number fails inside an accessor rather than with a custom message. A call evaluates the function and the argument (eager), then evaluates the body in \code{env} extended with \code{x=v2}: it is the \code{letC} clause with the body taken from the function value. The four \code{let} programs of \Cref{sec:envcheck} still give 37, 11, 12, 20, now as \code{numV}s, and simple functions work:
\begin{lstlisting}[style=repl]
> (eval (parse '((lambda x (+ 2 x)) 1)))
(numV 3)
> (eval (parse '(let x 2 (let addx (lambda y (+ y x)) (addx 3)))))
(numV 5)
\end{lstlisting}
``But now\ldots'' \code{env} in the \code{appC} clause is the environment of the \emph{call}. If the caller has rebound a name that the body uses freely, the body sees the caller's binding:
\begin{lstlisting}[style=repl]
(let x 2
  (let addx (lambda y (+ y x))
    (let x 5
      (addx 3))))
  env at the call:   [x=5, addx=(lambdaV y (+ y x)), x=2]
  body (+ y x) in    [y=3, x=5, addx=..., x=2]
    y ---> 3,  x ---> 5   (the caller's x)
---> 8
\end{lstlisting}
Lexically, the \code{x} in \code{addx}'s body is the \code{x=2} around the \code{lambda}, so the answer should be 5 (\Cref{fig:addx}). ``Oops, we have the same problem of dynamic scope!'' And for the same reason as \code{eval-env1}: the function body is stored code, and it is run in the environment where it is used.

\begin{figure}[htbp]
\centering
\begin{tikzpicture}[tok/.style={inner sep=0pt, font=\ttfamily\normalsize, anchor=base west}]
  \node[tok] at (0,0) {(let x 2 (let addx (lambda y (+ y x)) (let x 5 (addx 3))))};
  \node[tok] (ox) at (5\cw,0) {\phantom{x}};
  \node[tok] (ux) at (34\cw,0) {\phantom{x}};
  \node[tok] (ix) at (43\cw,0) {\phantom{x}};
  \draw[bind] (ux.south) to[out=-120, in=-60] (ox.south);
  \draw[bad] (ux.north) to[out=60, in=120] (ix.north);
  \node[lbl, text=accent] at (19\cw,-1.9) {lexical (\code{lecture5e}): the \code{x} around the \code{lambda}, so 5};
  \node[lbl, text=sigc] at (38.5\cw,0.8) {dynamic (\code{lecture5d}): the caller's \code{x}, so 8};
\end{tikzpicture}
\caption{The free \code{x} in \code{addx}'s body. Without closures, the body is evaluated in the caller's environment and finds the inner \code{x}.}
\label{fig:addx}
\end{figure}

This is also the week~4 eager-capture bug in a new form. In \Cref{sec:eagerapp}, \code{(let f (lambda x (+ x y)) (let y 6 (f 1)))} gave 7 under capturing substitution, because the free \code{y} of \code{f} met the \code{y=6} at the call. \code{lecture5d}'s \code{eval} gives 7 for it too. Capture in a substitution interpreter and dynamic scope in an environment interpreter are the same mistake.

\subsection{Attempt 2: Closures}

The fix is the one of \Cref{sec:dynscope}: save the environment where the code was written, here where the \code{lambda} was evaluated, together with the code.

\begin{definition}[Closure]
A \emph{closure} is a function value consisting of the parameter, the body, and the environment in which the \code{lambda} expression was evaluated. Calling it evaluates the body in that saved environment, extended with the parameter.
\end{definition}
\code{lecture5e.rkt} replaces \code{lambdaV} by \code{closV} and changes two clauses of \code{eval-env}; everything else is \code{lecture5d.rkt} unchanged:
\begin{lstlisting}
(define-type Value
  [numV (n : number)]
  [closV (env : Env) (x : symbol) (e : Expr)])

    [lambdaC (x e) (closV env x e)]
    ...
    [appC (e1 e2)
          (let [(v1 (eval-env env e1))
                (v2 (eval-env env e2))]
            ; run the body in the closure's saved env, not the current one
            (eval-env (extend-env (bind (closV-x v1) v2) (closV-env v1)) (closV-e v1)))]
\end{lstlisting}
\code{lambdaC} captures the current \code{env} in the closure. \code{appC} extends \code{(closV-env v1)}, the saved environment, instead of \code{env}, the caller's. \tyname{Value} and \tyname{Env} now refer to each other, as \tyname{Binding} and \tyname{Env} did in \code{lecture5c.rkt}. (\code{val->expr} is gone: nothing is substituted any more, so values never have to turn back into expressions.)

The file reruns the \code{addx} example with the inner \code{x} changed to 6 (dynamic scope would give 9):
\begin{lstlisting}[style=repl]
(let x 2
  (let addx (lambda y (+ y x))
    (let x 6
      (addx 3))))
  evaluating the lambda in [x=2]  ---> addx = (closV [x=2] y (+ y x))
  env at the call:   [x=6, addx, x=2]          (not used for the body)
  body (+ y x) in    [y=3, x=2]                (the closure's env + y=3)
    y ---> 3,  x ---> 2
---> 5
\end{lstlisting}
``It gives 5, as we'd expect with lexical scope.'' The closure carries its own environment, shown in \Cref{fig:closure}. Whatever the caller binds, the body sees only the bindings that were in scope where the \code{lambda} was written, plus its parameter. The week~4 example \code{(let f (lambda x (+ x y)) (let y 6 (f 1)))} is now \code{No binding found}, as it should be.

\begin{figure}[htbp]
\centering
\begin{tikzpicture}[x=1cm, y=1cm]
  \node[cell, fill=tsk, minimum width=14mm] (cenv) at (0,0) {env};
  \node[cell, fill=tsk, minimum width=8mm, right=0pt of cenv] (cx) {y};
  \node[cell, fill=tsk, minimum width=22mm, right=0pt of cx] (ce) {(+ y x)};
  \node[lbl, above=1mm of cx] {\code{addx} = \code{closV}};
  \node[cell, fill=uspace, minimum width=12mm] (b2) at (0,-1.6) {x=2};
  \node[right=2mm of b2, lbl] (nil) {\code{[]}};
  \draw[ptr] (b2) -- (nil);
  \draw[ptr] (cenv.center) -- (b2.north);
  \node[lbl, left=2mm of b2, align=right, text=accent] {saved when the\\ \code{lambda} was evaluated};

  \node[cell, fill=kspace, minimum width=12mm] (k1) at (6.4,-1.6) {x=6};
  \node[cell, fill=greybg, minimum width=14mm, right=3mm of k1] (k2) {addx};
  \node[cell, fill=uspace, minimum width=12mm, right=3mm of k2] (k3) {x=2};
  \draw[ptr] (k1) -- (k2); \draw[ptr] (k2) -- (k3);
  \node[lbl, below=1mm of k2] {caller's env at \code{(addx 3)}: ignored for the body};

  \node[cell, fill=kspace, minimum width=12mm] (y3) at (0,-3.2) {y=3};
  \draw[ptr] (y3) -- (b2);
  \node[lbl, right=2mm of y3, align=left] {body env = saved env + parameter:\\ \code{x} is found as 2};
\end{tikzpicture}
\caption{The closure for \code{addx} and the environments of the call \code{(addx 3)}. The body is evaluated in the closure's environment extended with \code{y=3}; the caller's \code{x=6} is never consulted.}
\label{fig:closure}
\end{figure}

\begin{keypoint}[Closures]
A function value must remember the environment in which it was created. This is not an optimization or a detail of this interpreter: it is the definition of lexical scope for first-class functions, and it is why \code{make-adder} in \Cref{sec:hof} returns a function that still knows its \code{x} after \code{make-adder} has returned.
\end{keypoint}

\subsection{The Five Files Side by Side}

Each file is a small change to the previous one, alternating between a first attempt that gets scope wrong and a fix:
\begin{center}
\begin{tabular}{@{}llll@{}}
\toprule
file & what is added & what a binding stores & scope \\
\midrule
\code{lecture5a} & environments, eager \code{let}, numbers only & a number & lexical \\
\code{lecture5b} & lazy \code{let}, attempt 1 & an unevaluated expression & \textcolor{sigc}{dynamic} \\
\code{lecture5c} & lazy \code{let}, attempt 2 & an expression \emph{and its env} & lexical \\
\code{lecture5d} & \code{lambda}, attempt 1 & a value; functions are \code{lambdaV} & \textcolor{sigc}{dynamic} \\
\code{lecture5e} & \code{lambda}, attempt 2 & a value; functions are closures & lexical \\
\bottomrule
\end{tabular}
\end{center}
The two red rows fail for the same reason and are fixed the same way: some code is run in the environment where it is used instead of the environment where it was written. On the programs of this chapter:
\begin{center}
\footnotesize
\begin{tabular}{@{}lccccc@{}}
\toprule
& substitution & \code{eval1} & \code{eval2} & \code{eval} & \code{eval} \\
& (week 4) & (5b) & (5c) & (5d) & (5e) \\
\midrule
\code{(let x 5 (let x 6 (+ x x)))} & 12 & 12 & 12 & 12 & 12 \\
\code{(let x 1 (let y x (let x 2 (+ y y))))} & 2 & \textcolor{sigc}{4} & 2 & 2 & 2 \\
\code{(let x y (let y 2 (+ x x)))} & error & \textcolor{sigc}{4} & error & error & error \\
\code{(let x (+ y y) 10)} & error & 10 & 10 & error & error \\
\code{addx} example & 5 & --- & --- & \textcolor{sigc}{8} & 5 \\
\code{(let f (lambda x (+ x y)) (let y 6 (f 1)))} & error & --- & --- & \textcolor{sigc}{7} & error \\
\bottomrule
\end{tabular}
\end{center}
The substitution column is the capture-avoiding eager interpreter \code{eval-eager-v3} (\Cref{sec:eagerapp}). Red entries disagree with it for scope reasons; ``---'' means the file has no \code{lambda}. The row \code{(let x (+ y y) 10)} is not a scope difference but the eager/lazy one of \Cref{sec:lazy}: lazy interpreters never evaluate an unused \code{e1}.
